% Physical pages: 422--429 (printed pages 399--406).
% Figures: none.
% Uncertainties: none.

\subsection{Path Integrals for Fermions}

We now turn to the problem of extending the path-integral formalism to cover
theories containing fermions as well as bosons. It would be easy to proceed in
a purely formal way, by analogy with the bosonic case, with the justification
that this gives the `right' Feynman rules. Instead, we will here derive the
path-integral formalism for fermions directly from the principles of quantum
mechanics, as we did for bosons~\hyperref[ref:9.9]{[9]}.

As before, we will start with a general quantum mechanical system, with
`coordinates' $Q_a$ and canonical conjugate `momenta' $P_a$, but now satisfying
anticommutation rather than commutation relations:
\begin{equation}
\{Q_a,P_b\}=i\delta_{ab},
\tag{9.5.1}\label{eq:9.5.1}
\end{equation}
\begin{equation}
\{Q_a,Q_b\}=\{P_a,P_b\}=0.
\tag{9.5.2}\label{eq:9.5.2}
\end{equation}
(These are Schrödinger-picture operators, or in other words
Heisenberg-picture operators at time $t=0$.) Later we will replace the discrete
index $a$ with a spatial position $\mathbf{x}$ and a field index $m$.

We wish first to construct a complete basis for the states on which the $Q$s
and $P$s act. Note that for any given $a$, we have
\begin{equation}
Q_a^2=P_a^2=0.
\tag{9.5.3}\label{eq:9.5.3}
\end{equation}
It follows that there will always be a `ket' state $\ket{\Omega}$ annihilated
by all $Q_a$:
\begin{equation}
Q_a\ket{\Omega}=0,
\tag{9.5.4}\label{eq:9.5.4}
\end{equation}
and a `bra' state $\bra{\Omega}$ annihilated (from the right) by all $P_a$:
\begin{equation}
\bra{\Omega}P_a=0.
\tag{9.5.5}\label{eq:9.5.5}
\end{equation}
For instance, we can take
\[
\ket{\Omega}\propto\left(\prod_a Q_a\right)\ket{f},
\qquad
\bra{\Omega}\propto\bra{g}\left(\prod_a P_a\right),
\]
where $\ket{f}$ and $\bra{g}$ are any kets and bras for which these expressions
do not vanish. (They cannot vanish for all $\ket{f}$ and $\bra{g}$ unless the
operators $\prod_aQ_a$ and $\prod_aP_a$ vanish, which we assume not to be the
case.) These states satisfy Eqs. \eqref{eq:9.5.4} and \eqref{eq:9.5.5} by virtue of Eq. \eqref{eq:9.5.3}.
They are not in general unique, because there may be other bosonic degrees of
freedom that distinguish the various possible $\ket{\Omega}$ and
$\bra{\Omega}$, but for simplicity we will limit ourselves here to the case
where the only degrees of freedom are those described by the fermionic
operators $Q_a$ and $P_a$, and will assume that the states satisfying Eqs.
\eqref{eq:9.5.4} and \eqref{eq:9.5.5} are unique up to constant factors, which we choose so that
\begin{equation}
\braket{\Omega}{\Omega}=1.
\tag{9.5.6}\label{eq:9.5.6}
\end{equation}
(Note that this normalization convention could not be imposed if we had
defined $\bra{\Omega}$ as the left-eigenstate of the $Q_a$ with eigenvalue
zero, because in this case $\bra{\Omega}\{Q_a,P_b\}\ket{\Omega}$ would vanish,
which with Eq. \eqref{eq:9.5.1} would imply that $\braket{\Omega}{\Omega}=0$.)

As we saw in Section 7.5, in the Dirac theory $Q_a$ is not Hermitian, but
instead has an adjoint $-iP_a$, in which case $\bra{\Omega}$ can be regarded as
simply the adjoint of $\ket{\Omega}$. However, there are fermionic operators
(such as the `ghost' fields to be introduced in Volume II) for which $P_a$ is
unrelated to the adjoint of $Q_a$. In what follows we will not need to assume
anything about the adjoints of $Q_a$ or $P_a$, or about any relation between
$\ket{\Omega}$ and $\bra{\Omega}$.

A complete basis for the states of this system is provided by $\ket{\Omega}$
and the states (antisymmetric in indices $a,b,\ldots$)
\begin{equation}
\ket{a,b,\ldots}\equiv P_aP_b\cdots\ket{\Omega}
\tag{9.5.7}\label{eq:9.5.7}
\end{equation}
with any number of \emph{different} $P$s acting on $\ket{\Omega}$. That is, the
result of acting on these states with any operator function of the $P$s and
$Q$s can be written as a linear combination of the same set of states. In
particular, if an index $a$ is unequal to any of the indices appearing in
$\ket{b,c,\ldots}$, then
\begin{equation}
Q_a\ket{b,c,\ldots}=0,
\tag{9.5.8}\label{eq:9.5.8}
\end{equation}
\begin{equation}
P_a\ket{b,c,\ldots}=\ket{a,b,c,\ldots}.
\tag{9.5.9}\label{eq:9.5.9}
\end{equation}
On the other hand, if $a$ is equal to one of the indices in the sequence,
$b,c,\ldots$, we can always rewrite the state (possibly changing its sign) so
that $a$ is the first of these indices, in which case we have
\begin{equation}
Q_a\ket{a,b,c,\ldots}=i\ket{b,c,\ldots},
\tag{9.5.10}\label{eq:9.5.10}
\end{equation}
\begin{equation}
P_a\ket{a,b,c,\ldots}=0.
\tag{9.5.11}\label{eq:9.5.11}
\end{equation}

Similarly, we may define a complete dual basis, consisting of $\bra{\Omega}$
and the states (also antisymmetric in the indices)
\begin{equation}
\bra{a,b,\ldots}\equiv\bra{\Omega}\cdots(-iQ_b)(-iQ_a).
\tag{9.5.12}\label{eq:9.5.12}
\end{equation}
Using Eqs. \eqref{eq:9.5.4}--\eqref{eq:9.5.6} and the anticommutation relation \eqref{eq:9.5.1}, we see
that the scalar products of these states take the values
\begin{align}
\braket{c,d,\ldots}{a,b,\ldots}
&=\bra{\Omega}\cdots(-iQ_d)(-iQ_c)P_aP_b\cdots\ket{\Omega}\notag\\
&=\begin{cases}
0&\text{if }\{c,d,\ldots\}\ne\{a,b,\ldots\},\\
1&\text{if }c=a,\ d=b,\ \text{etc.}
\end{cases}
\tag{9.5.13}\label{eq:9.5.13}
\end{align}
where $\{\cdots\}$ here denotes the set of indices within the brackets,
irrespective of order.

In deriving the Feynman rules, we would like to be able to rewrite sums over
intermediate states like \eqref{eq:9.5.7} as integrals over eigenstates of the $Q_a$ or
the $P_a$. However, it is not possible for these operators to have eigenvalues
(other than zero) in the usual sense. Suppose we try to find a state $\ket{q}$
that satisfies (for all $a$)
\begin{equation}
Q_a\ket{q}=q_a\ket{q}.
\tag{9.5.14}\label{eq:9.5.14}
\end{equation}
From Eq. \eqref{eq:9.5.2} we see that
\begin{equation}
q_aq_b+q_bq_a=0,
\tag{9.5.15}\label{eq:9.5.15}
\end{equation}
which is impossible for ordinary numbers. However, nothing can stop us from
introducing an algebra of `variables' (known as \emph{Grassmann variables})
$q_a$, which act like c-numbers as far as the physical Hilbert space is
concerned, but which still satisfy the anticommutation relations \eqref{eq:9.5.15}. We
will require further that
\begin{equation}
\{q_a,q'_b\}=\{q_a,Q_b\}=\{q_a,P_b\}=0,
\tag{9.5.16}\label{eq:9.5.16}
\end{equation}
where $q$ and $q'$ denote any two `values' of these variables. We can now
construct eigenstates $\ket{q}$ satisfying Eq. \eqref{eq:9.5.14}:
\begin{equation}
\ket{q}=\exp\left(-i\sum_aP_aq_a\right)\ket{\Omega}
\tag{9.5.17}\label{eq:9.5.17}
\end{equation}
with the exponential defined as usual by its power series expansion. (To
verify Eq. \eqref{eq:9.5.14}, use the fact that all $P_aq_a$ commute with one another
and have zero square, so that
\begin{align*}
[Q_a-q_a]\ket{q}
&=[Q_a-q_a]\exp(-iP_aq_a)
 \exp\left(-i\sum_{b\ne a}P_bq_b\right)\ket{\Omega}\\
&=[Q_a-q_a][1-iP_aq_a]
 \exp\left(-i\sum_{b\ne a}P_bq_b\right)\ket{\Omega}\\
&=[-i\{Q_a,P_a\}q_a-q_a]
 \exp\left(-i\sum_{b\ne a}P_bq_b\right)\ket{\Omega}=0
\end{align*}
as required by Eq. \eqref{eq:9.5.14}.) We can also define left-eigenstates $\bra{q}$
(\emph{not} the adjoints of $\ket{q}$), as
\begin{align}
\bra{q}
&\equiv\bra{\Omega}\left(\prod_aQ_a\right)
 \exp\left(-i\sum_aq_aP_a\right)\notag\\
&=\bra{\Omega}\left(\prod_aQ_a\right)
 \exp\left(+i\sum_aP_aq_a\right),
\tag{9.5.18}\label{eq:9.5.18}
\end{align}
where $\prod_a$ is the product in whatever order we take as standard. By the
same argument as for Eq. \eqref{eq:9.5.14}, we see that
\begin{equation}
\bra{q}Q_a=\bra{q}q_a.
\tag{9.5.19}\label{eq:9.5.19}
\end{equation}

These eigenstates have the scalar product
\begin{align*}
\braket{q'}{q}
&=\bra{\Omega}\left(\prod_aQ_a\right)
 \exp\left(i\sum_bP_b(q'_b-q_b)\right)\ket{\Omega}\\
&=\bra{\Omega}\left(\prod_aQ_a\right)
 \left(\prod_b[1+iP_b(q'_b-q_b)]\right)\ket{\Omega}.
\end{align*}
Moving each $Q_a$ to the right (starting with the rightmost) yields factors
$i^2(q'_a-q_a)$, which we move to the right out of the scalar product, so
\begin{equation}
\braket{q'}{q}=\prod_a(q_a-q'_a).
\tag{9.5.20}\label{eq:9.5.20}
\end{equation}
We shall see that Eq. \eqref{eq:9.5.20} plays the role of a delta function in integrals
over the $q$s.

In the same way, we can construct right- and left-eigenstates of the $P_a$:
\begin{equation}
P_a\ket{p}=p_a\ket{p},
\tag{9.5.21}\label{eq:9.5.21}
\end{equation}
\begin{equation}
\bra{p}P_a=\bra{p}p_a,
\tag{9.5.22}\label{eq:9.5.22}
\end{equation}
where the $p_a$ are like $q_a$ anticommuting c-numbers (taken for convenience
to anticommute with the $q_a$ and all fermionic operators as well as each
other), and
\begin{equation}
\ket{p}=\exp\left(-i\sum_aQ_ap_a\right)
 \left(\prod_bP_b\right)\ket{\Omega},
\tag{9.5.23}\label{eq:9.5.23}
\end{equation}
\begin{equation}
\bra{p}=\bra{\Omega}\exp\left(-i\sum_a p_aQ_a\right)
\tag{9.5.24}\label{eq:9.5.24}
\end{equation}
with scalar product (now derived by moving the $P$s to the left)
\begin{equation}
\braket{p'}{p}=\prod_a(p'_a-p_a).
\tag{9.5.25}\label{eq:9.5.25}
\end{equation}
The scalar products of these two sorts of eigenstate with each other are
\begin{align*}
\braket{q}{p}
&=\bra{q}\exp\left(-i\sum_aQ_ap_a\right)
 \left(\prod_aP_a\right)\ket{\Omega}\\
&=\left(\prod_a\exp(-iq_ap_a)\right)
 \bra{q}\left(\prod_aP_a\right)\ket{\Omega}\\
&=\left(\prod_a\exp(-iq_ap_a)\right)
 \bra{\Omega}\left(\prod_aQ_a\right)
 \left(\prod_aP_a\right)\ket{\Omega},
\end{align*}
and so
\begin{equation}
\braket{q}{p}=\chi_N\exp\left(-i\sum_aq_ap_a\right)
=\chi_N\exp\left(i\sum_ap_aq_a\right),
\tag{9.5.26}\label{eq:9.5.26}
\end{equation}
where $\chi_N$ is a phase that depends only on the number $N$ of $Q_a$
operators:
\[
\chi_N\equiv\bra{\Omega}\left(\prod_aQ_a\right)
 \left(\prod_aP_a\right)\ket{\Omega}
=i^N(-1)^{N(N-1)/2}.
\]
Somewhat more simply, we also find
\begin{equation}
\braket{p}{q}=\prod_a\exp(-ip_aq_a).
\tag{9.5.27}\label{eq:9.5.27}
\end{equation}

It is easy to see that the states $\ket{q}$ are in a sense a complete set (and
so also are the $\ket{p}$.) From the definitions \eqref{eq:9.5.17}, we see that the
state $\ket{a,b,\ldots}$ in the general basis is (up to a phase) just the
coefficient of the product $q_aq_b\cdots$ in an expansion of $\ket{q}$ in a
sum of products of $q$s. Therefore we can write any state $\ket{f}$ in the
form
\[
\ket{f}=f_0\ket{q}_0+\sum_a f_a\ket{q}_a
 +\sum_{a\ne b}f_{ab}\ket{q}_{ab}+\cdots,
\]
where the $f$s are numerical coefficients, and a subscript $a,b,\ldots$ on
$\ket{q}$ denotes the coefficient of $q_aq_b\cdots$ in $\ket{q}$.

In summing over states, it will be very convenient to introduce a sort of
integration over fermionic variables, known as \emph{Berezin integration}~\hyperref[ref:9.10]{[10]},
that is designed to pick out the coefficients of such products of
anticommuting c-numbers. For any set of such variables $\xi_n$ (either $p$s or
$q$s or both together), the most general function $f(\xi)$ (either a c-number
or a state-vector like $\ket{q}$) can be put in the form
\begin{equation}
f(\xi)=\left(\prod_n\xi_n\right)c
 +\text{terms with fewer $\xi$ factors}
\tag{9.5.28}\label{eq:9.5.28}
\end{equation}
and the integral over the $\xi$s is defined simply by
\begin{equation}
\int\left(\widetilde{\prod_n}d\xi_n\right)f(\xi)\equiv c
\tag{9.5.29}\label{eq:9.5.29}
\end{equation}
with the tilde in Eq. \eqref{eq:9.5.29} indicating that we use the convenient convention
that the differentials are written in an order \emph{opposite} to that of the
product of integration variables in Eq. \eqref{eq:9.5.28}. Since this product is
antisymmetric under the interchange of any two $\xi$s, the integral is
likewise antisymmetric under the interchange of any two $d\xi$s, so these
`differentials' effectively anticommute
\begin{equation}
d\xi_n\,d\xi_m+d\xi_m\,d\xi_n=0.
\tag{9.5.30}\label{eq:9.5.30}
\end{equation}
Also, the coefficient $c$ may itself depend on other unintegrated c-number
variables that anticommute with the $\xi$s over which we integrate, in which
case it is important to standardize the definition of $c$ by moving all $\xi$s
to the left of $c$ before integrating over them, as we have done in Eq.
\eqref{eq:9.5.28}.

For instance, the most general function of a pair of anticommuting c-numbers
$\xi_1$ and $\xi_2$ takes the form
\[
f(\xi_1,\xi_2)=\xi_1\xi_2c_{12}+\xi_1c_1+\xi_2c_2+d
\]
because the squares and all higher powers of $\xi_1$ and $\xi_2$ vanish. This
function has the integrals
\[
\int d\xi_1\,f(\xi_1,\xi_2)=\xi_2c_{12}+c_1,
\qquad
\int d\xi_2\,f(\xi_1,\xi_2)=-\xi_1c_{12}+c_2,
\]
\[
\int d\xi_2\,d\xi_1\,f(\xi_1,\xi_2)=c_{12}.
\]
Note that the multiple integral is the same as a repeated integral:
\[
\int d\xi_2\,d\xi_1\,f(\xi_1,\xi_2)
=\int d\xi_2\left[\int d\xi_1\,f(\xi_1,\xi_2)\right],
\]
a result that can easily be extended to integrals over any number of fermionic
variables. (It was in order to obtain this result without extra sign factors
that we took the product of differentials in Eq. \eqref{eq:9.5.29} to be in the opposite
order to the product of variables in Eq. \eqref{eq:9.5.28}.) Indeed, we could have first
defined the integral over a single anticommuting c-number $\xi_1$, and then
defined multiple integrals in the usual way by iteration. The most general
function of anticommuting c-numbers is linear in any one of them
\[
f(\xi_1,\xi_2,\ldots)=b(\xi_2,\ldots)+\xi_1c(\xi_2,\ldots)
\]
(because $\xi_1^2=0$), and its integral over $\xi_1$ is defined as
\[
\int d\xi_1\,f(\xi_1,\xi_2,\ldots)=c(\xi_2,\ldots).
\]
Repeating this process leads to the same multiple integral as defined by Eqs.
\eqref{eq:9.5.28} and \eqref{eq:9.5.29}.

This definition of integration shares some other properties with multiple
integrals (from $-\infty$ to $+\infty$) over ordinary real variables, but there
are significant differences.

Obviously, Berezin integration is linear, in the sense that
\begin{align}
\int\left(\widetilde{\prod_n}d\xi_n\right)[f(\xi)+g(\xi)]
={}&\int\left(\widetilde{\prod_n}d\xi_n\right)f(\xi)
+\int\left(\widetilde{\prod_n}d\xi_n\right)g(\xi)
\tag{9.5.31}\label{eq:9.5.31}
\end{align}
and also
\begin{equation}
\int\left(\widetilde{\prod_n}d\xi_n\right)[f(\xi)a(\xi')]
=\left[\int\left(\widetilde{\prod_n}d\xi_n\right)f(\xi)\right]a(\xi'),
\tag{9.5.32}\label{eq:9.5.32}
\end{equation}
where $a(\xi')$ is any function (including a constant) of any anticommuting
c-numbers $\xi'_m$ over which we are \emph{not} integrating. However, linearity
with respect to left-multiplication is not so obvious. If we are integrating
over $\nu$ variables, then since $\xi'_m$ is assumed to anticommute with all
$\xi_n$, we have
\[
a((-1)^\nu\xi')\left(\prod_n\xi_n\right)
=\left(\prod_n\xi_n\right)a(\xi')
\]
and so
\begin{equation}
\int\left(\widetilde{\prod_n}d\xi_n\right)
[a((-1)^\nu\xi')f(\xi)]
=a(\xi')\int\left(\widetilde{\prod_n}d\xi_n\right)f(\xi).
\tag{9.5.33}\label{eq:9.5.33}
\end{equation}
It is therefore very convenient (though not strictly necessary) to take the
differentials $d\xi_n$ to anticommute with all anticommuting variables
(including the $\xi_n$):
\begin{equation}
(d\xi_n)\xi'_m+\xi'_m(d\xi_n)=0
\tag{9.5.34}\label{eq:9.5.34}
\end{equation}
in which case Eq. \eqref{eq:9.5.33} reads more simply
\begin{equation}
\int a(\xi')\left(\widetilde{\prod_n}d\xi_n\right)f(\xi)
=a(\xi')\int\left(\widetilde{\prod_n}d\xi_n\right)f(\xi).
\tag{9.5.35}\label{eq:9.5.35}
\end{equation}
Another similarity with ordinary integration is that, for an arbitrary
anticommuting c-number $\xi'$ independent of $\xi$,
\begin{equation}
\int\left(\widetilde{\prod_n}d\xi_n\right)f(\xi+\xi')
=\int\left(\widetilde{\prod_n}d\xi_n\right)f(\xi)
\tag{9.5.36}\label{eq:9.5.36}
\end{equation}
since shifting $\xi$ by a constant only affects the terms in $f$ with fewer
than the total number of $\xi$-variables.

On the other hand, consider a change of variables
\begin{equation}
\xi_n\longrightarrow\xi'_n=\sum_m\mathcal{G}_{nm}\xi_m,
\tag{9.5.37}\label{eq:9.5.37}
\end{equation}
where $\mathcal{G}$ is an arbitrary non-singular matrix of ordinary numbers.
The product of the new variables is
\[
\prod_n\xi'_n
=\sum_{m_1m_2\cdots}\left(\prod_n\mathcal{G}_{nm_n}\xi_{m_n}\right).
\]
But $\prod_n\xi_{m_n}$ here is just the same as the product (in the original
order) $\prod_n\xi_n$, except for a sign $\epsilon[m]$ which is $+1$ or $-1$
according to whether the permutation $n\to m_n$ is an even or odd permutation
of the original order:
\[
\prod_n\xi'_n
=\left[\sum_{m_1m_2\cdots}
 \left(\prod_n\mathcal{G}_{nm_n}\right)\epsilon[m]\right]
 \prod_n\xi_n
=(\operatorname{Det}\mathcal{G})\prod_n\xi_n.
\]
This applies whatever order we take for the $\xi_n$, as long as we take the
$\xi'_n$ in the same order. It follows that the coefficient of
$\prod_n\xi'_n$ in any function $f(\xi)$ is just
$(\operatorname{Det}\mathcal{G})^{-1}$ times the coefficient of
$\prod_n\xi_n$, a statement we write as
\begin{equation}
\int\left(\widetilde{\prod_n}d\xi'_n\right)f
=(\operatorname{Det}\mathcal{G})^{-1}
 \int\left(\widetilde{\prod_n}d\xi_n\right)f.
\tag{9.5.38}\label{eq:9.5.38}
\end{equation}
This is the usual rule for changing variables of integration, except that
$(\operatorname{Det}\mathcal{G})$ appears to the power $-1$ instead of $+1$.
We shall use Eq. \eqref{eq:9.5.38} and the linearity properties \eqref{eq:9.5.31}, \eqref{eq:9.5.32}, and
\eqref{eq:9.5.35} later to evaluate the integrals encountered in deriving the Feynman
rules for theories with fermions.

We can now use this definition of integration to write the completeness
condition as a formula for an integral over eigenvalues. As already mentioned,
any state $\ket{f}$ can be expanded in a series of the states $\ket{\Omega}$,
$\ket{a}$, $\ket{a,b}$, etc. and these states are (up to a phase) the
coefficients of the products $1$, $q_a$, $q_aq_b$, etc. in the $Q$-eigenstate
$\ket{q}$. According to the definition of integration here, we can pick out
the coefficient of any product $q_bq_cq_d\cdots$ in the state $\ket{q}$ by
integrating the product of $\ket{q}$ with all $q_a$ with $a$ \emph{not} equal
to $b,c,d,\ldots$. Thus, by choosing a function $f(q)$ as a suitable sum of
such products of $q$s, we can write any state $\ket{f}$ as an integral:
\begin{equation}
\ket{f}=\int\left(\widetilde{\prod_a}dq_a\right)\ket{q}f(q)
=\int\ket{q}\left(\widetilde{\prod_a}dq_a\right)f(q).
\tag{9.5.39}\label{eq:9.5.39}
\end{equation}
(We can move $\ket{q}$ to the left of the differentials without any sign
changes because the exponential in Eq. \eqref{eq:9.5.17} used to define $\ket{q}$
involves only even numbers of fermionic quantities.) To find the function
$f(q)$ for a given state-vector $\ket{f}$, take the scalar product of Eq.
\eqref{eq:9.5.39} with some bra $\bra{q'}$ (with $q'$ any fixed $Q$-eigenvalue).
According to Eqs. \eqref{eq:9.5.35} and \eqref{eq:9.5.20}, this is
\[
\braket{q'}{f}=\int\left(\prod_a(q_a-q'_a)\right)
 \left(\widetilde{\prod_b}dq_b\right)f(q).
\]
Moving every factor $(q_a-q'_a)$ to the right past every differential $dq_b$
yields a sign factor $(-)^{N^2}=(-)^N$, where $N$ is now the total number of
$q_a$ variables, so
\[
\braket{q'}{f}=(-)^N\int\left(\widetilde{\prod_b}dq_b\right)
 \left(\prod_a(q_a-q'_a)\right)f(q).
\]
We can rewrite $f(q)$ as $f(q'+(q-q'))$ and expand in powers of $q-q'$. All
terms beyond the lowest order vanish when multiplied with the product
